\documentclass[11pt]{article}

\usepackage[margin=0.95in]{geometry}
\usepackage[T1]{fontenc}
\usepackage[utf8]{inputenc}
\usepackage{lmodern}
\usepackage{microtype}
\usepackage[hidelinks]{hyperref}
\usepackage{booktabs}
\usepackage{xcolor}
\usepackage{amsmath}
\usepackage{array}

\definecolor{accent}{RGB}{18,54,92}
\setlength{\parindent}{0pt}
\setlength{\parskip}{8pt}

\title{\textbf{A Full Explanation of the QQQ Systematic Signal Research Project}}
\author{}
\date{\today}

\begin{document}
\maketitle
\thispagestyle{empty}

\section*{Introduction}

This project began with a simple but serious question: can a disciplined set of
interpretable timing rules improve the risk-adjusted performance of QQQ relative
to simply buying it and holding it? The attraction of QQQ is obvious. It gives
concentrated exposure to large-cap growth and technology, and over long periods
it has delivered strong upside. The difficulty is equally obvious. QQQ also
comes with deep drawdowns, sharp reversals, and periods in which trend,
volatility, and macro conditions interact in ways that make passive exposure
painful. The project therefore does not ask whether QQQ is a good asset in the
long run. It asks whether exposure to QQQ can be managed more intelligently.

The central idea is not stock picking and it is not black-box prediction. The
project is about \emph{exposure management}. On each day, the system decides how
much QQQ exposure should be held on the next day. Sometimes that desired
exposure is strongly long. Sometimes it is flat. Sometimes it is slightly
short. Those decisions are driven by a small set of understandable signals
related to trend, oversold conditions, volatility regimes, cross-asset
confirmation, and seasonality. The project is therefore best understood as a
market-timing research pipeline with explicit backtesting, validation, and
report generation, rather than as a one-off trading notebook.

\section*{What the Project Is Trying to Solve}

Buy-and-hold QQQ has a structural weakness: it treats all market states the
same. A strong uptrend, a volatility shock, an oversold crash, and a weak
macro regime all receive the same portfolio response. That is attractive for
simplicity, but it means the investor absorbs the full force of every drawdown
and every regime transition. The project studies whether a diversified set of
rules can respond differently to those environments without collapsing into
curve-fitted complexity.

The research problem is therefore not ``can I predict tomorrow's QQQ return
perfectly?'' That would be both unrealistic and conceptually vague. The real
problem is narrower and more defensible. If QQQ is above trend, perhaps it
deserves more exposure. If QQQ is violently oversold, perhaps a mean-reversion
signal should dominate. If volatility becomes extreme, perhaps gross risk should
be cut. If the dollar is strengthening while QQQ trend deteriorates, perhaps a
macro confirmation filter should turn more defensive. The project asks whether a
well-designed combination of such signals can improve the return-risk tradeoff
once one accounts for transaction costs and out-of-sample validation.

\section*{The Data and the Public-versus-Private Distinction}

One of the most important things to understand about the repository is that it
contains both a public, reproducible research framework and a documented summary
of the original Quanta QR Fellowship finalist submission. These are related, but
they are not identical. The repository ships with a sanitized sample
dataset spanning \texttt{2018-01-02} through \texttt{2025-06-30}. At minimum,
that dataset contains the columns \texttt{date}, \texttt{qqq\_close},
\texttt{qqq\_high}, \texttt{qqq\_low}, and \texttt{dxy\_close}. Those columns
are enough to reproduce the signal mechanics, backtest logic, result tables, and
plots without redistributing private challenge data.

The original challenge project, however, used a longer private history. Its
intended split was:

\begin{center}
\begin{tabular}{lll}
\toprule
\textbf{Split} & \textbf{Dates} & \textbf{Purpose} \\
\midrule
Train & 2000-01-01 to 2015-12-31 & Hypothesis development \\
Validation & 2016-01-01 to 2021-12-31 & Signal selection and ensemble design \\
Blind holdout & 2022-01-01 to 2025-06-30 & Final out-of-sample evaluation \\
\bottomrule
\end{tabular}
\end{center}

Because the included sample begins in 2018, the committed results adapts the
split in a transparent way:

\begin{center}
\begin{tabular}{lll}
\toprule
\textbf{Public split} & \textbf{Dates} & \textbf{Purpose} \\
\midrule
Train & 2018-01-01 to 2020-12-31 & First-pass design and filtering \\
Validation & 2021-01-01 to 2022-12-31 & Selection of final sleeves \\
Holdout & 2023-01-01 to 2025-06-30 & Final public out-of-sample evaluation \\
\bottomrule
\end{tabular}
\end{center}

This distinction matters because it explains why the public results are not
supposed to reproduce the original headline Sharpe. The included sample exists so
that the mechanics of the project are inspectable. The original challenge
summary exists so that the stronger private-run result is not erased. The
repository is explicit about this split because hiding it would make the project
less credible, not more.

\section*{What a Signal Means in This Project}

Each signal in the project outputs a daily target exposure to QQQ. Exposure is
bounded in the range \([-1.0, 1.5]\), where \(-1.0\) means fully short QQQ,
\(0.0\) means no exposure, \(1.0\) means fully long QQQ, and \(1.5\) means
levered long exposure. This framing is important because it clarifies what the
system is doing. The project is not trying to classify each day as ``up'' or
``down.'' It is deciding how aggressively or defensively to hold QQQ given the
information currently available.

The signal set was chosen to remain interpretable. The trend sleeves ask whether
QQQ is above long or medium moving averages. The volatility and mean-reversion
sleeves ask whether realized range-based volatility is low, high, or extreme,
and whether RSI indicates a deeply oversold condition. The risk-control sleeves
look for shock behavior or skew patterns that may make short exposure dangerous.
The macro sleeve uses DXY trend as a confirmation filter. The seasonality sleeve
encodes turn-of-month behavior. None of these are intended to be mystical. The
point is to build a diversified zoo of transparent hypotheses and test them
under a common framework.

\section*{How the Individual Signals Behave}

The simplest signal in the repo is the long-term trend sleeve. If QQQ is above a
long moving average, the model holds a high long exposure. If QQQ is below that
long-term trend threshold, it goes to cash. This is a classic trend-following
rule and serves as a baseline for the idea that persistent strength should be
respected rather than faded.

The medium-term trend sleeve is more tactical. It uses a shorter exponential
moving average and allows mild short exposure when QQQ drops below that trend.
This makes it more responsive and more active than the long-term sleeve, which
is why it becomes the strongest trend representative on the included sample.

The volatility and mean-reversion logic is concentrated in \texttt{rsi\_deep\_value}.
This sleeve is designed around the idea that deeply oversold selloffs should not
always be treated as invitations to go risk-off. If volatility is elevated but
the market is extremely oversold, the sleeve overrides the usual defensive logic
and becomes aggressively long, attempting to capture rebound behavior rather than
selling the hole.

The defensive sleeves, such as \texttt{rsi\_gated\_short},
\texttt{conservative\_fade}, \texttt{vol\_shock\_dampener}, and
\texttt{skew\_filter}, are all attempts to formalize the same broad intuition:
there are environments in which naive shorting or naive persistence of trend
exposure becomes dangerous. Those sleeves cut or neutralize exposure when
volatility, range, or skew patterns become too extreme.

The macro sleeve, \texttt{dual\_trend\_macro}, asks whether QQQ trend should be
trusted equally when the dollar is strengthening versus when it is weakening. It
is a simple way of encoding the idea that macro risk appetite may not be fully
captured by QQQ price alone. Finally, the \texttt{turn\_of\_month} sleeve captures
a simple calendar effect by raising exposure around month-end and the start of a
new month.

\section*{How the Backtest Works}

The backtest is intentionally conservative. A signal is computed at time \(t\),
but the resulting exposure is applied to the \emph{next} close-to-close return,
not the same-day return. That one-day shift is an important design choice,
because it avoids accidental look-ahead bias. The project is therefore trying to
answer the question, ``if I knew today's information, what exposure should I
hold on the next day?''

For each strategy, the backtest records exposure, applied exposure, gross
return, turnover, transaction cost, and net return. Transaction costs are
modeled as one-way cost per unit turnover, with the repo using a
default assumption of 5 basis points. Turnover is measured directly from changes
in exposure from one day to the next. This means that more active signals are
penalized appropriately, and any apparently attractive sleeve has to survive not
only out-of-sample evaluation but also trading friction.

The resulting backtest is therefore not merely a curve of cumulative returns. It
is a mechanically explicit mapping from signal series to implemented next-day
exposures, and from those exposures to cost-adjusted P\&L.

\section*{How the Final Ensemble Was Chosen}

The final public ensemble was not formed by averaging every signal in the zoo.
Instead, the project uses a validation-only ranking process. Each signal is
assigned to a family. Within each family, sleeves are ranked using a score that
rewards higher validation Sharpe and shallower drawdowns while penalizing
turnover and annualized cost drag. Signals with near-duplicate validation
profiles are filtered out so that the final ensemble does not simply repeat the
same behavior under multiple names.

On the committed included sample, that process leaves three sleeves, which can be
described in plain language as a medium-term trend sleeve, an oversold
mean-reversion sleeve, and a turn-of-month sleeve. The final public ensemble is
the clipped combination of those three. This is conceptually useful because it means the final blend
contains one trend sleeve, one oversold/volatility sleeve, and one seasonality
sleeve. In other words, the ensemble is trying to diversify \emph{research
logic}, not just average many versions of the same idea.

\section*{How to Read the Main Metrics}

Sharpe ratio measures return per unit of volatility, so it is the main
risk-adjusted performance metric in the project. Max drawdown measures the worst
peak-to-trough loss. Turnover measures how aggressively exposure changes from day
to day. Cost drag translates that turnover into performance lost to friction.
Validation metrics are used for selection, while holdout metrics are used for
final out-of-sample evaluation. This distinction is crucial. A good-looking
validation result is only interesting if it survives on holdout after the
selection process is already fixed.

\section*{What the Public Results Actually Show}

The headline public results at 5 bps one-way cost are:

\begin{center}
\small
\begin{tabular}{@{}>{\raggedright\arraybackslash}p{0.28\textwidth} >{\raggedleft\arraybackslash}p{0.10\textwidth} >{\raggedleft\arraybackslash}p{0.12\textwidth} >{\raggedleft\arraybackslash}p{0.13\textwidth} >{\raggedleft\arraybackslash}p{0.12\textwidth} >{\raggedleft\arraybackslash}p{0.09\textwidth}@{}}
\toprule
\textbf{Strategy} & \shortstack{\textbf{Val}\\\textbf{Sharpe}} & \shortstack{\textbf{Holdout}\\\textbf{Sharpe}} & \shortstack{\textbf{Holdout}\\\textbf{Return}} & \shortstack{\textbf{Holdout}\\\textbf{Max DD}} & \shortstack{\textbf{Turn}\\\textbf{over}} \\
\midrule
Buy \& Hold QQQ & 0.26 & -0.45 & -26.0\% & -40.1\% & 0.00 \\
Medium-Term Trend & 0.65 & 0.02 & -2.9\% & -21.0\% & 0.09 \\
RSI Deep Value & 0.26 & -0.45 & -39.0\% & -54.9\% & 0.00 \\
Dual Trend Macro & -0.34 & -0.91 & -40.2\% & -44.6\% & 0.10 \\
Final Research Ensemble & 0.36 & -0.29 & -17.0\% & -31.9\% & 0.06 \\
\bottomrule
\end{tabular}
\normalsize
\end{center}

The most important thing about this table is that it is \emph{not flattering},
and that is part of what makes it credible. The public-sample final ensemble did
not produce positive holdout Sharpe. The medium-term trend sleeve is the
cleanest individual public performer, but even that sleeve has only slightly
positive holdout Sharpe. The final research ensemble, however, did improve
drawdown versus buy-and-hold: its holdout max drawdown was materially shallower
than the passive benchmark. In other words, the included sample suggests that
exposure management helped more on downside control than on outright alpha.

That matters conceptually. A serious research project is allowed to conclude
that a included sample is not strong enough to support a pretty headline. In fact,
keeping that result in the repo is evidence that the project prioritizes
research discipline over cosmetic storytelling.

\section*{What the Original Quanta Challenge Result Was}

The original Quanta QR Fellowship finalist submission is documented separately
in the original challenge summary report because the longer private dataset is
not redistributed. The preserved challenge materials report that the final
ensemble achieved net blind-holdout Sharpe above \(1.3\) after costs over the
period from January 1, 2022 through June 30, 2025. That is the result behind
the CV line.

So the honest interpretation of the project is the following. The repo
demonstrates the mechanics: signal construction, no-lookahead backtesting,
validation-only selection, cost modeling, sensitivity analysis, regime
diagnostics, and report generation. The original private-run summary is what
supports the stronger headline performance claim. The repository is careful not
to blur those two things together.

\section*{What the Project Lets Us Infer}

The strongest inference is not that the included sample proves a durable edge. The
strongest inference is that a disciplined timing research process can be built
around simple and interpretable sleeves, and that drawdown-aware exposure
management can be meaningfully separated from pure buy-and-hold behavior. The
project also shows that family diversity matters. Several defensive sleeves
turned out to be near-duplicates of \texttt{rsi\_deep\_value} on the public
sample, and that finding is just as informative as a strong positive metric
would have been. It tells us that naming a signal differently is not the same as
having a genuinely distinct source of behavior.

The project also demonstrates that validation and holdout discipline are not
minor implementation details. An earlier version of the report generation path
had holdout leakage in the selection logic. That was removed, and the final
public ensemble is now selected with a validation-only score. This is one of the
most important research decisions in the whole repo, because it reinforces the
principle that out-of-sample results should be reported, not optimized.

\section*{Limitations}

The repository is careful about its own limits. The public data is
sanitized and shorter than the intended full-history design. DXY and
cross-asset behavior are only as realistic as the committed sample permits.
Daily close-to-close execution is idealized. Transaction costs are simplified.
And the stronger original challenge result is still documented as a preserved
summary rather than a fully reproducible raw historical backtest. None of these
limitations destroy the value of the project, but they do define the right
interpretation. This is a research framework and documented research record,
not a production trading system.

\section*{Conclusion}

The QQQ project is best understood as a serious systematic research pipeline for
timing exposure to a single highly liquid ETF using a diversified set of simple
signals. Its real strength is methodological clarity: interpretable sleeves,
explicit train / validation / holdout structure, no-lookahead mechanics,
transaction-cost accounting, committed results, and honest retention of
unflattering public-sample results. That is precisely why it works well as a
 project. It does not ask the reader to take the process on
faith. It shows the process directly, and then clearly distinguishes that public
framework from the stronger private-run result that originally motivated it.

\end{document}
